\section{Problem Statement}
\label{sec:problem}

\subsection{Statement}
\textit{Using only the measurements from the first $N\le100$
charge--discharge cycles of a lithium-ion cell, predict its cycle life
(the number of cycles until capacity falls to 80\% of nominal).
Systematically compare machine-learning models from different families and
data-driven against physics-informed feature sets, under one common and
leakage-free evaluation protocol. Assess how well the best models carry
over to an unseen batch tested under different conditions.}

\subsection{Research questions}
\begin{enumerate}
  \item[\textbf{RQ1}] Can the early-prediction results of~\cite{severson2019}
  be reproduced on its public dataset?
  \item[\textbf{RQ2}] With about 124 cells, do other model types (whole-curve
  PLS, kernel methods, additive splines) outperform regularised linear models?
  \item[\textbf{RQ3}] Do physics-informed features (differential capacity,
  resistance growth, heating) improve on the purely data-driven features of
  \cite{severson2019}, and for which models?
  \item[\textbf{RQ4}] How does prediction error depend on the number of early
  cycles $N$ that are used?
  \item[\textbf{RQ5}] How well do the trained models carry over to an unseen
  batch with a different test procedure (2018-02-20)?
\end{enumerate}

\subsection{Data}
\begin{itemize}
  \item \textbf{Development data:} the three public batches
  of~\cite{severson2019,matr_data} (2017-05-12, 2017-06-30, 2018-04-12).
  These contain 124 LFP/graphite cells with lifetimes from 150 to 2300
  cycles. We use the original split into 41 training, 43 primary-test and
  40 secondary-test cells, so results are directly comparable with the
  paper.
  \item \textbf{Generalisation test:} the 2018-02-20 batch analysed in this
  report (39 labelled cells, Section~\ref{sec:dataset}). It is never used
  for training or tuning, and errors are reported separately for standard
  and \texttt{-newstructure} cells.
\end{itemize}

\subsection{Feature sets}
All features are computed only from cycles $1\ldots N$ (default $N=100$).
Table~\ref{tab:featsets} lists the four feature sets.

\begin{table}[t]
\centering
\caption{Feature sets compared in this study.}
\label{tab:featsets}
\setlength{\tabcolsep}{3pt}
\begin{tabular}{@{}lp{0.70\columnwidth}@{}}
\toprule
Set & Features \\
\midrule
F1 Variance & $\log_{10}\mathrm{Var}(\dq)$ (1 feature)~\cite{severson2019} \\
F2 Discharge & F1 + min, skewness, kurtosis of $\dq$; discharge capacity at
cycle 2; maximum capacity minus capacity at cycle 2 (6)~\cite{severson2019} \\
F3 Full & min and $\log_{10}$Var of $\dq$; slope and intercept of a linear
fit to capacity fade over cycles 2--100; discharge capacity at cycle 2;
average charge time over the first 5 cycles; integral of temperature over
cycles 2--100; minimum internal resistance; change in internal resistance
from cycle 2 to cycle 100 (9)~\cite{severson2019} \\
F4 Physics & capacity margin to EOL over the $\dq$ amplitude; early
capacity rise after cycle 10 (overhang) (2); selected from 19 physics-based
candidates (Section~\ref{sec:physics}) \\
\bottomrule
\end{tabular}
\end{table}

\subsection{Models}
Table~\ref{tab:models} lists the models. The original plan of seven models
was reduced to five: tree ensembles and neural networks were dropped because
with about 41 training cells they can memorise the training set, and trees
cannot predict lives outside the training range (Section~\ref{sec:models}).

\begin{table}[t]
\centering
\caption{Models compared in this study.}
\label{tab:models}
\setlength{\tabcolsep}{3pt}
\begin{tabular}{@{}llp{0.42\columnwidth}@{}}
\toprule
Family & Model & Reason for inclusion \\
\midrule
\multirow{2}{*}{Linear} & Elastic net & Baseline of~\cite{severson2019}; easy to interpret \\
 & PLS regression & Built for many correlated inputs; uses the whole $\dq$ curve \\
\midrule
\multirow{2}{*}{Kernel} & SVR & Non-linear; robust loss \\
 & Gaussian process & Non-linear, and gives prediction uncertainty \\
\midrule
Additive & Spline model (GAM) & One smooth, readable curve per feature \\
\bottomrule
\end{tabular}
\end{table}

\subsection{Evaluation protocol}
\begin{enumerate}
  \item \textbf{Target:} $\log_{10}$(cycle life). Predictions are converted
  back to cycles before computing errors.
  \item \textbf{Splits:} the same split for every model. Hyperparameters are
  tuned only by cross-validation on the training set. The test sets are
  used once, for the final evaluation.
  \item \textbf{Metrics:} RMSE (cycles), mean absolute percentage error
  (MAPE, \%) and $R^2$, reported separately for each test set.
  \item \textbf{Baselines:} predict the training mean, and predict the mean
  for the cell's charging policy. A model has to beat both before its
  result is counted as meaningful.
  \item \textbf{Repeats:} randomised models are run with 5--10 seeds and
  reported as mean $\pm$ standard deviation.
\end{enumerate}

\subsection{Additional analyses}
\begin{itemize}
  \item \textbf{Early-prediction horizon (RQ4):} recompute features for
  $N\in\{20,30,40,50,60,80,100,150\}$ and plot test error against $N$
  (Section~\ref{sec:models}).
  \item \textbf{Interpretability:} use permutation importance and SHAP values
  to find which features, and therefore which degradation mechanisms,
  drive the predictions.
\end{itemize}

\subsection{Scope}
\textbf{In scope:} predicting cycle life from early cycles; comparing
models and feature sets; testing generalisation across batches.
\textbf{Out of scope:} predicting the full capacity trajectory,
physics-informed neural networks (PINNs), electrochemical simulation, and
degradation-mode analysis that needs reference electrode (half-cell) data.

\subsection{Expected deliverables}
\begin{enumerate}
  \item A reproducible pipeline from the raw \texttt{.mat} files to features
  and models (code in the project repository).
  \item A results table (models $\times$ feature sets $\times$ test sets).
  \item Plots of predicted against actual cycle life, error against $N$,
  and feature importance.
  \item An analysis of generalisation to the unseen batch.
\end{enumerate}
